\section{Fonctionnalités attendues}

% ----------------------------------------------------------------
% 📌 À quoi sert cette section ?
% ----------------------------------------------------------------
% Cette partie permet de décrire les actions concrètes que le produit final permettra d'accomplir.
% Il s’agit de traduire les objectifs fonctionnels en fonctionnalités mesurables.
% Cela permet aussi de poser les bases de l’évaluation du projet.

% ----------------------------------------------------------------
% 🧭 Comment la remplir ?
% ----------------------------------------------------------------
% ➤ Décrivez les fonctionnalités sous forme d’User Stories :
%     - "En tant que [utilisateur], je veux [fonction] afin de [bénéfice]."
% ➤ Estimez le temps nécessaire pour chaque fonctionnalité (en JH = jour-homme)
% ➤ Organisez les fonctionnalités par modules, groupes logiques ou priorités si pertinent.
% ➤ Vous pouvez utiliser une liste ou un tableau synthétique.

% ----------------------------------------------------------------
% 🧾 Exemple de liste d’User Stories :
% ----------------------------------------------------------------
% \begin{itemize}
%   \item \textbf{[1 JH]} En tant qu’étudiant, je veux coder du SVG avec aperçu en temps réel.
%   \item \textbf{[2 JH]} En tant qu’utilisateur, je veux valider des défis à difficulté croissante.
%   \item \textbf{[1 JH]} En tant que formateur, je veux accéder à des statistiques de progression.
% \end{itemize}

% ----------------------------------------------------------------
% 📊 Exemple de tableau synthétique :
% ----------------------------------------------------------------
% \begin{table}[H]
% \centering
% \begin{tabular}{|p{8cm}|c|}
% \hline
% Fonctionnalité (User Story) & Estimation (JH) \\
% \hline
% Codage SVG avec aperçu en temps réel & 1 \\
% Défis de reproduction graphique gamifiés & 2 \\
% Interface stats/formateur & 1 \\
% Documentation interactive & 1 \\
% \hline
% \textbf{Total estimé} & \textbf{5} \\
% \hline
% \end{tabular}
% \caption{Fonctionnalités attendues et estimation de charge}
% \end{table}

% ----------------------------------------------------------------
% ✍️ À vous de jouer : complétez ci-dessous
% ----------------------------------------------------------------

% (Suggestion de départ)
Les fonctionnalités prévues sont :

\begin{itemize}
  \item [À compléter...]
\end{itemize}
